

\section{Proposed Activities}\label{sec:activities}

We have categorized the workshop goals into two broad themes, and we will have workshop sessions with a dedicated theme. The first $75$ mins session of the workshop will be dedicated to the theme {\em "Communication and Education of Uncertainty"} ($\mathbf{T_1}$). This session will include workshop activities related to the topics of uncertainty propagation, communication, and education. The second 75 mins session of the workshop will be dedicated to the theme {\em "AI and Decision-Making"} ($\mathbf{T_2}$), which will include workshop activities focused on AI uncertainty and decision-making topics. 
We will divide each session into two types of activities, specifically paper presentations and a panel session.  Table~\ref{workshopAgenda} summarizes the time budget for each of the proposed activities. 

\begin{table}[!h]
\centering
\caption{Workshop Agenda}
\label{workshopAgenda}
%
\resizebox{0.75\linewidth}{!}{
\begin{tabular}{lr}
%\toprule
%Activity & Time Budget \\
\toprule
\multicolumn{2}{c}{\textbf{$\mathbf{T_1}$: Communication and Education of Uncertainty}} \\
\midrule
Introduction &  5 mins \\
Paper Presentations with Q\&A & 25 mins \\
Panel & 35 mins \\
Panel Q\&A & 10 mins \\
\midrule
Break & 30 mins \\
\midrule
\multicolumn{2}{c}{\textbf{$\mathbf{T_2}$: AI and Decision-Making}} \\
\midrule
Paper Presentations with Q\&A & 25 mins \\
Panel & 35 mins \\
Panel Q\&A & 10 mins \\
Closing & 5 mins \\
\bottomrule
\end{tabular}}
\vspace{-3mm}
\end{table}


% we will have each $75$ mins session of the workshop with a dedicated theme.

{\large Presentations:} The 2024 uncertainty workshop received $12$ total submissions, among which $6$ were accepted for short/long presentation. We, therefore, again anticipate $6$-$8$ accepted submissions for presentation at the workshop this year. Each presenting author will have around 5-7 minutes to present their work depending upon the number of accepted submissions as well as if the submitted paper is short (4 pages) or long (9 pages). We currently anticipate $3$-$4$ papers for each of the themes $\mathbf{T_1}$ and $\mathbf{T_2}$. However, in the case of uneven distribution of accepted papers across the two themes, we will let the audience know about this skewed distribution upfront in the introduction. The audience will have 2 minutes to ask questions after each presentation. In the case that there is no enough time to ask questions, the audience will have the opportunity to bring up important questions/topics/issues during the  $10$ mins Q\&A in the panel session. In addition, the workshop participants will be encouraged to comment and ask questions on a discord channel (or whatever platform is used by IEEE VIS) to have opportunities to discuss those during the panel session or informal coffee breaks. 

{\large Panel:} We intend to do one panel session for each of the two themes $\mathbf{T_1}$ and $\mathbf{T_2}$. We plan to invite $3$ panelists for each panel session, thereby providing on average $10$ mins to each panelist to share there persepectives on various issues. For the theme $\mathbf{T_1}$, we plan to invite one application scientist as panelist, preferably from outside the visualization domain, who can shed light on the critical nature of uncertainty propagation and communication in their domain. In addition, we will invite one expert in uncertainty visualization who can share views on essentials (e.g., confidence intervals, probability) that non-experts should be educated with to enable interpretation of uncertainty. Lastly, we will invite expert in the visualization field to share views on intersection of visualization, application, and uncertainty.  

For the panel session for theme $\mathbf{T_2}$, we plan to invite two experts in the field of AI who can discuss various critical aspects, including but not limited to sources of uncertainty in AI models, approaches to enable trustworthy scientific discovery using AI, and role of AI in decision-making under uncertainty. We also intend to invite an expert in the psychology field who can share their perspectives on the role of perception and cognition in decision-making under uncertainty. The panel members will discuss and debate over questions asked by the session moderator and audience. Each session will be moderated by one of the workshop organizers to ensure the smooth flow of the session in both physical or virtual setting. The one request for the proposed panel session would be multiple microphones. 

%As such, the only financial support we request is free conference attendee registration and support for their travel expenses for any keynote or invited speakers (if possible).


%The total time budget of the proposed activities is 2 hours and 40 mins. This leaves 20 mins for a coffee break in between the proposed activities.



%{\large Keynote:} $\;$ Our intention is to invite renowned keynote speakers who will discuss challenges and potential solutions relevant to each of the mentioned themes ($\mathbf{T_1}$ and $\mathbf{T_2}$). We will stress more on Q\&A session at the end of each eynote talk to foster interaction and gain understanding of future directions. For the theme $\mathbf{T_1}$, we plan to invite application scientist, preferably from outside the visualization domain, who can shed light on critical nature of uncertainty propagation and communication in their domain. For the theme $\mathbf{T_2}$, we plan to invite expert in the field of AI or psychology who can discuss critical factors relevant to decision-making when interpreting results with uncertainty. As such, the only financial support we request is free conference attendee registration and support for their travel expenses for any keynote or invited speakers (if possible).

%{\large Posters:} The authors of the accepted posters will present their preliminary ideas in uncertainty analysis in visualization or their application domains. Given time constrains of the workshop, we will request the VIS chairs to have posters be presented during the main VIS poster event. However, we propose a lightening talk round by the authors of the accepted posters during the workshop to provide a summary and encourage the audience to attend their poster presentations. We currently anticipate 3-5 accepted posters for this workshop. Posters is expected to be a productive interactive session where the authors will get some interesting feedback from applications scientists with diverse backgrounds and the opportunity to engage in future collaborations. 










% Our workshop will mainly comprise four activities, a keynote presentation, posters, paper presentations, and a panel session.

% {\large Keynote:} $\;$ Our intention is to invite a renowned keynote speaker, preferably application scientist outside the visualization domain, who can discuss why uncertainty should be considered the first-class citizen of the data analysis through use cases in their domain. As such, the only financial support we request is free conference attendee registration and support for their travel expenses for any keynote or invited speakers.

% {\large Posters:} The authors of the accepted posters will present their preliminary ideas in uncertainty analysis in visualization or their application domains. Given time constrains of the workshop, we will request the VIS chairs to have posters be presented during the main VIS poster event. However, we propose a lightening talk round by the authors of the accepted posters during the workshop to provide a summary and encourage the audience to attend their poster presentations. We currently anticipate 3-5 accepted posters for this workshop. Posters is expected to be a productive interactive session where the authors will get some interesting feedback from applications scientists with diverse backgrounds and the opportunity to engage in future collaborations. 

% %The workshop organizers will also set up a couple of demos (e.g., isosurface uncertainty using VTK-m~\cite{TA:2023:Wang:funmc2} in Figure~\ref{fig:teaser}c) along with the authors of the accepted submissions. For the live demos, we will prefer arrangement of bigger screens, but these could be also presented on laptops/desktops. We currently anticipate 5 poster slots/software demos in the final workshop.

% {\large Presentations:} The authors of the accepted papers will present their work
% in around 5-7 minutes depending upon the number of accepted submissions. We currently anticipate 3-5 accepted submissions for presentation at the workshop. If the number of accepted papers is more than anticipated, we will balance it with the panel (discussed next) time. The audience will get 2 minutes to ask questions. In the case, there is no sufficient time to ask questions, the audience will get the opportunity to bring up important questions/topics/issues in the panel session. Additionally, the workshop participants will be encouraged to
% comment and ask questions on a discord (or whatever platform is
% used by IEEE VIS) channel that could be discussed during the panel session or informal coffee breaks. 

% {\large Panel:} We intend to invite five experienced scientists with expertise spanned across data science, visualization, and applications to share their perspectives on uncertainty and interact with the audience. In particular, our intended theme is to conduct a panel session on the topic of {\em "Top Challenges in Uncertainty Visualization and Analysis"}. The panel members will discuss and debate over questions asked by the audience. The session will be moderated by the workshop organizers, but it will be mainly intended to promote interaction on the uncertainty topic. Another reason for having the session moderated is to ensure the smoothness of a session in both physical or virtual setting. The one request for the proposed panel session would be multiple microphones. 


% The Table~\ref{workshopAgenda} summarizes the time budget for each of the proposed activities. The total time budget of the proposed activities is 2 hours and 40 mins. This leaves 20 mins for a coffee break in between the proposed activities.
% \begin{table}[h]
% \centering
% \caption{Workshop Agenda}
% \label{workshopAgenda}
% %
% \begin{tabular}{lll}
% \toprule
% Activity & Time Budget \\
% \midrule
% Introduction &  5 mins \\
% Keynote & 45 mins \\
% Keynote Q\&A & 10 mins \\
% Poster lightning talks & 15 mins \\
% \midrule
% Break & 20 mins \\
% \midrule
% Paper Presentations & 40 mins \\
% Panel & 40 mins \\
% Closing & 5 mins \\
% \bottomrule
% \end{tabular}
% \vspace{-3mm}
% \end{table}

% %\begin{table}[h]
% %\centering
% %\caption{Workshop Agenda}
% %\label{workshopAgenda}
% %%
% %\begin{tabular}{lll}
% %\toprule
% %Activity & Presenter & Time Budget \\
% %\midrule
% %Introduction & Athawale & 5 mins \\
% %Keynote & XYZ & 45 mins \\
% %Keynote Q\&A & XYZ & 10 mins \\
% %Poster lightning talks & 3-5 presenters & 15 mins \\
% %Break & & 30 mins \\
% %Paper Presentations & 3-5 Presenters & 30 mins \\
% %Panel & XYZ & 40 mins \\
% %Closing & Athawale & 5 mins \\
% %\bottomrule
% %\end{tabular}
% %\end{table}

% %\ta{Panel:}\\
% %- set up a panel on "top challenges in uncertainty vis"\\
% %- Jessica Hullman, Matt Kay\\
% %- Uncertainty visualization for machine learning (Prasanna Balkrishna at ORNL)\\
% %- Maybe a few other domain experts  who can talk about uncertainty (Dan Lu in climate science team at ORNL)\\\\
% %\ta{Posters, live demos:}\\
% %- We will invite live demos of software usage in addition to posters and presentations for enhancing understanding of existing software tools for uncertainty visualization.\\
% %- plugins have been developed that enable visualization of uncertainty in isosurfaces using the VTK-m library and visualization of uncertainty in critical points using the topology toolkit.